\chapter{La città come parte dello strumento}\label{sec:city}

Una città può ospitare un'altra città nello spazio di un incontro musicale.
Una persona si siede al Portale di Bergamo e, davanti a sé, incontra un'altra
persona seduta a un secondo Portale. Lo schermo OLED e la videocamera
cercano una presenza quasi speculare: riconoscere il volto, vedere le mani,
intuire un'esitazione, cominciare insieme. La disposizione e la scala
visiva devono essere provate; la reciprocità è già un criterio del progetto.
Entrambe le persone devono poter vedere, ascoltare, proporre e interrompere.

Il nucleo dell'esperienza è suonare contemporaneamente attraverso due
luoghi. La connessione rende percepibile una distanza che per un momento
si può attraversare con un gesto musicale. La qualità tecnica del percorso
fra i due nodi decide quali gesti e repertori possono realmente convivere:
la misura della rete deve precedere la promessa pubblica del concerto.
Quando quella qualità non basta, una conversazione o uno scambio a turni
possono conservare valore, ma vanno annunciati come modalità diverse.

La bellezza civica di questa proposta risiede nella simmetria dell'invito.
Il partner remoto non fornisce un contenuto da consumare: ospita persone
che hanno a loro volta una città, un luogo e un pubblico. Curare questo
incontro significa dare a entrambi i lati tempo e iniziativa. Il Portale
può diventare una piccola infrastruttura di ospitalità reciproca, nella
quale la musica precede talvolta le parole e aiuta a trovarle.

\section{Caprotti: una scelta da argomentare}
Parco Caprotti appartiene alla geografia del documento originale del 2016;
il Padiglione del Tè è invece una candidatura sviluppata successivamente.
Questa distinzione impedisce di attribuire alla fonte un progetto localizzato
che essa non documenta. L'interesse per il Padiglione nasce dall'ipotesi
che una presenza musicale osservabile dal giardino possa costruire una
continuità fra attraversamento, curiosità e partecipazione. È un'ipotesi
spaziale da confrontare con un rilievo, con gli usi e con le condizioni reali
di accesso: trasparenza del vetro e disponibilità di uno spazio non sono
equivalenti.

Due tracce documentarie rendono più precisa la storia civica del parco.
L'inventario dell'Ufficio Tecnico registra nel 1947 una manifestazione di
volontà di donazione al Comune: documento 1435, faldone 16, pagina PDF 44.
Il documento originale non è stato consultato; l'inventario non prova
l'avvenuto trasferimento. Alla pagina PDF 65 compaiono il documento 2075,
associato a Luigi Angelini e datato 1948, e il 2076, relativo alla
Soprintendenza e datato 1950, nel faldone 30. Sono piste per una ricerca
sulle trasformazioni del giardino, non attribuzioni del progetto originario
del Padiglione \cite{mai_inventario}. La raccolta comunale IBCAA riproduce,
alla pagina PDF 176, il vincolo 118 sul Giardino Caprotti e zona di rispetto,
con decreto del 12 agosto 1949 e notifica del 10 febbraio 1950
\cite{bergamo_vincolo118}. La scheda non risolve da sola la disciplina applicabile
a un nuovo allestimento.

La verifica storica successiva dovrà ricercare progetto originario,
trasformazioni e atti relativi al manufatto. Il fascicolo di fattibilità
dovrà separatamente acquisire titolo e condizioni d'uso attuali,
competenza degli interlocutori e percorso delle eventuali autorizzazioni.
Queste due ricerche rispondono a domande diverse. Una ricostruisce ciò
che è accaduto; l'altra stabilisce ciò che possiamo proporre oggi.

\section{Un luogo già abitato da pratiche}
La pagina ufficiale del Cerchio di Gesso dedicata al \emph{Giardino
dell'anima} descrive attività di recupero e cura delle piante con volontari
e annuncia una Festa Tanabata il 7 luglio 2026 alla Sala del té di Parco
Caprotti \cite{cerchio2026}. È prova della programmazione pubblicata,
non dell'effettivo svolgimento dell'evento, della disponibilità quotidiana
della sala o dei diritti sullo spazio. È però sufficiente per richiedere
un confronto con persone e attività già presenti.

La proposta potrebbe trovare una relazione fertile con pratiche di cura:
tempi lenti, attenzione alle differenze, gesti che acquistano valore quando
vengono ripetuti insieme. Questa è una lettura progettuale, non una
collaborazione già concordata. Non autorizza a sonorizzare automaticamente
le attività esistenti. L'incontro preliminare dovrebbe produrre una mappa
condivisa degli orari, degli ingombri, dei depositi, dei momenti che
richiedono silenzio e delle attività eventualmente compatibili. Il primo
risultato utile potrebbe essere una sola giornata comune, oppure la
decisione motivata che questo luogo non sia adatto.

\section{Una coppia di luoghi, una relazione da costruire}
L'unità di progetto per gli incontri remoti è una coppia di Portali.
Un nodo può ospitare studio individuale o attività locali, ma un collegamento
civico richiede una seconda organizzazione capace di accogliere le persone,
curare il programma e garantire il presidio. Nessuna città, biblioteca,
scuola o associazione è qui indicata come partner acquisito.
La selezione dovrebbe considerare interesse reciproco, accesso pubblico,
competenze di accompagnamento, disponibilità simultanea e compatibilità
tecnica, oltre alla distanza geografica.

Si possono confrontare tre assetti: due presenze temporanee per una serie
limitata di incontri; una coppia stabile con calendario regolare; un nodo
continuativo che incontra a rotazione nodi temporanei compatibili.
Il primo assetto riduce l'impegno sul sito ma moltiplica montaggi e trasporti.
Il secondo sostiene il ritorno e le relazioni, ma richiede risorse durevoli
in entrambi i luoghi. Il terzo amplia gli interlocutori e richiede una nuova
verifica di rete, audio, immagine e gestione a ogni accoppiamento.
La compatibilità di un nodo non dimostra quella di tutte le sue connessioni.

La cura del programma può iniziare con due musicisti che si presentano
attraverso una breve frase, proseguire con variazioni condivise e aprirsi
alla partecipazione dei presenti. Musiche, età e lingue vanno scelte con
entrambi gli ospitanti. Si può lasciare a ciascuna città la proposta di una
parte degli incontri, alternando chi invita e chi risponde. Il gemellaggio
musicale è qui una possibilità progettuale, non uno status istituzionale
attribuito ai luoghi.

Fra gli appuntamenti, la modalità studio mantiene vivo il nodo: esercizi,
materiali e accompagnamento possono provenire da una piattaforma propria
o da un servizio esterno selezionato. La scelta dovrà confrontare
accessibilità, funzionamento senza account personale ove possibile,
licenze, esportazione dei materiali e dipendenza dal fornitore. Uno
strumento didattico esterno non va presentato come tecnologia originale
del Portale; la sua integrazione e la relazione con gli incontri pubblici
sono invece parte del progetto. L'eventuale piattaforma propria richiede
un budget distinto e un ambito funzionale definito.

La collocazione di ciascun nodo deve superare le verifiche su disponibilità,
accessibilità, compatibilità acustica, tutela del luogo e reversibilità.
Le evidenze mancanti restano tali: un punteggio estetico favorevole non può
compensare l'assenza di un titolo d'uso o di un percorso di accesso praticabile.

\section{Valore pubblico, attrattività e diritto di passare oltre}
L'incontro con una persona lontana può essere una porta d'ingresso, purché non diventi l'unico
criterio. Un oggetto sorprendente può generare fotografie senza produrre
scambio musicale. Il progetto dovrebbe osservare separatamente curiosità
iniziale, possibilità di cominciare, collaborazione fra persone e ritorno
nel tempo. Una visita breve può avere valore; non dimostra apprendimento.
Un numero elevato di presenze non dimostra accesso equo.

La promessa civica comprende anche chi non suona. Il giardino deve poter
continuare a ospitare riposo, conversazioni e attraversamento. Per questo
la programmazione dovrà stabilire modalità sonore, fasce temporali,
procedura di segnalazione e responsabile capace di interrompere una
sessione. Il dissenso di un residente o di un frequentatore è un dato di
progetto da esaminare, non un ostacolo comunicativo da neutralizzare.

L'accesso equo richiede scelte osservabili: ingresso gratuito alle prove
pubbliche di base; avvio senza account personale; istruzioni comprensibili;
possibilità di osservare senza partecipare; turni che non favoriscano
sempre chi ha già esperienza. Vanno provati avvicinamento, postura,
raggiungibilità dei comandi e forme alternative di interazione con persone
che presentano esigenze diverse. La gratuità non elimina barriere motorie,
sensoriali, linguistiche o sociali. Prima di parlare di accessibilità
realizzata servono prove e correzioni documentate.


La finestra verso l'altro luogo deve avere un'apertura dichiarata e una
chiusura riconoscibile. Chi attraversa il giardino non deve diventare
involontariamente comparsa di una trasmissione. Inquadratura, segnalazione
della connessione e area di partecipazione vanno progettate insieme.
La proposta assume incontri concordati e presidiati; non una finestra
permanentemente aperta sui frequentatori. Il pubblico deve sapere se vede
una persona in diretta, un materiale registrato o una schermata di studio.
